% GDZ PPN236018647,native538lower–540above55;HKUST539–541.
\NumberedItem{50}
\[
\left.\begin{aligned}
\int\surd\sin x.dx&=2\int\frac{yy\,dy}{\surd(1-y^4)}\\[.5ex]
\int\surd\operatorname{tang}x.dx&=2\int\frac{dy}{{}^{4}\!\surd(1-y^4)}\\[.5ex]
\int\surd\frac1{\sin x}.dx&=2\int\frac{dy}{\surd(1-y^4)}
\end{aligned}\right\}\quad yy=\genfrac{}{}{0pt}{}{\sin}{\cos}x.
\]
\noindent\hfill1797 Ian. 7.

\medskip\begingroup\small
\indent Zu dem ersten und dritten der hier auftretenden Integrale vergl. die in dem Abschnitt [I.] der \textit{Ältesten Untersuchungen über lemniskatische Funktionen}, art. [1.], S. 145 und art. [3.], S. 146 wiedergegebenen
\par\endgroup
\clearpage\SourceLeaf{539}\ArithmeticHead{510}{L}{TAGEBUCH MIT ERLÄUTERUNGEN.}
\begingroup\small
\noindent Leisteaufzeichnungen, sowie die dazugehörigen Bemerkungen S. 149, wo auch Literaturnachweise gegeben sind. Das zweite Integral ist ein besonderer Fall des in der Nr. 53 vom 12. Januar 1797 aufgezeichneten Integrals \(\displaystyle\int\frac{dx}{\sqrt[n]{1-x^n}}\), das durch die Substitution\footnote{Für \(n=4\) ist diese Substitution bei S. 27 des \textsc{Leiste} ausgeführt; ebenda bei S. 104 finden sich mit der Überschrift »\textls[120]{Untersuchungen über das Integral} \(\displaystyle y=\int\frac{dx}{{}^{4}\!\surd(1-x^4)}\)« Reihenentwicklungen und numerische Rechnungen.}\ \(z^n=\frac{1-x^n}{x^n}\) in \(\displaystyle-\int\frac{z^{n-2}\,dz}{1+z^n}\) übergeht und so wieder als besonderer Fall des in der Nr. 54 betrachteten Integrals \(\displaystyle\int\frac{x^n\,dx}{1+x^m}\) erscheint.
\par\noindent\hfill\textsc{Schlesinger.}
\par\endgroup

\NumberedItem{51}
\indent Curvam (elasticam) lemniscatam a
\[
\int\frac{dx}{\surd(1-x^4)}
\]
\noindent pendentem perscrutari coepi.
\par\noindent\hfill[1797] Ian. 8.

\medskip\begingroup\small
\indent Damit stimmt überein die auf dem letzten Blatte der Scheda Ac befindliche Notiz: »Functiones lemniscaticas considerare coeperamus 1797. Januar. 8.«, siehe oben S. 206. Daß \textsc{Gauss} schon früher die Umkehrung des lemniskatischen Integrals in eine Reihe entwickelt hatte (siehe den art. [6.] der \textit{Exercitationes mathematicae}, oben S. 141), steht mit diesen Angaben nicht im Widerspruch, sondern zeigt, daß \textsc{Gauss} erst am 8. Januar 1797 die volle Bedeutung des lemniskatischen Integrals erkannt und es zum Ausgangspunkte einer selbständigen Theorie gemacht hat, während es ihm früher nur als ein einzelner Fall des allgemeinen Integrals
\[
\int\frac{dx}{\surd(1-x^n)}
\]
\noindent erschienen war (siehe die Tagebuchaufzeichnungen Nr. 32 und Nr. 33 vom September 1796). Die wahre Eigenart des Falles \(n=4\) hat \textsc{Gauss} durch Anwendung des \textsc{Euler}schen Additionstheorems erkannt; dies zeigen die oben S. 147 ff. in den artt. [4.]—[7.] abgedruckten Leistenotizen. Daß diese Notizen mit den am 8. Januar 1797 begonnenen Untersuchungen in Verbindung stehen, geht auch schon aus dem äußerlichen Merkmal hervor, daß hier wie dort in der Überschrift das Wort \textit{elastica} durchstrichen und an seine Stelle \textit{lemniscata} gesetzt ist.
\par\noindent\hfill\textsc{Klein. Schlesinger.}
\par\endgroup

\NumberedItem{52}
\indent Criterii \textsc{Euler}iani rationem sponte detexi.
\par\noindent\hfill[1797] Ian. 10.

\medskip\begingroup\small
\indent Die beiden folgenden Notizen Nr. 53, 54 (vergl. auch Nr. 50) zeigen, daß es sich um das \textsc{Euler}sche (richtiger \textsc{Newton}sche) Kriterium für die sogenannten binomischen Integrale \(\int x^{m-1}(a+bx^n)^{\frac\mu\nu}\,dx\) handelt. Siehe \textsc{Euler}, \textit{Institutiones calculi integralis} I, 1768, § 104, Opera omnia, ser. I, vol. 11, S. 62.
\par\noindent\hfill\textsc{Schlesinger.}
\par\endgroup

\clearpage\SourceLeaf{540}\ArithmeticHead{511}{R}{1797 IAN. 7—1797 IAN. 19.}
\NumberedItem{53}
\indent Integrale completum
\[
\int\frac{dx}{{}^{n}\!\surd(1-x^n)}
\]
\noindent ad circ[uli] quadr[aturam] reducere commentus sum.
\par\noindent\hfill[1797] Ian. 12.

\medskip\begingroup\small
\indent Vergl. die Bemerkung zu der Nr. 50 und die Nr. 59 vom 2. März 1797. Im § 352 von \textsc{Eulers} \textit{Institutiones calculi integralis} I, 1768, Opera omnia, ser. I, vol. 11, S. 226, wird für dieses \textit{integrale completum} der Wert
\[
\int\limits_0^1\frac{dx}{{}^{n}\!\surd(1-x^n)}=\frac\pi{n\sin\frac\pi n}
\]
\noindent gegeben; im § 353 werden dann die einfachsten Fälle \(n=2,3,4,6\) berechnet. \textsc{Gauss} will sagen, daß, da \(\sin\frac\pi n\) für ein beliebiges \(n\) durch Wurzeln aus ganzen Zahlen dargestellt werden kann, der betreffende Integralwert, was seinen transzendenten Charakter angeht, nur von \(\pi\) abhängt.
\par\noindent\hfill\textsc{Schlesinger.}
\par\endgroup

\NumberedItem{54}
\indent Methodus facilis
\[
\int\frac{x^n\,dx}{1+x^m}
\]
\noindent determinandi.
\par\noindent\hfill[1797 Ian.]

\medskip\begingroup\small
\indent Hierher gehörige Formeln und Rechnungen findet man aufgezeichnet im \textsc{Leiste} bei den Seiten 27, 91, 92; sie schließen sich durchweg an die Methode an, nach der \textsc{Euler} diese Integrale behandelt, siehe \textit{Institutiones calculi integralis} I, 1768, art. 77., Opera omnia, ser. I, vol. 11, S. 41, oder auch \textit{Methodus integrandi formulas differentiales rationales unicam variabilem involventes}, Comm. Acad. Petrop. 14 (1744/6), 1751, S. 3, Opera omnia, ser. I, vol. 17, S. 70, besonders die artt. 44—59.
\par\noindent\hfill\textsc{Klein. Schlesinger.}
\par\endgroup
